%2multibyte Version: 5.50.0.2960 CodePage: 1252

\documentclass{article}
%%%%%%%%%%%%%%%%%%%%%%%%%%%%%%%%%%%%%%%%%%%%%%%%%%%%%%%%%%%%%%%%%%%%%%%%%%%%%%%%%%%%%%%%%%%%%%%%%%%%%%%%%%%%%%%%%%%%%%%%%%%%%%%%%%%%%%%%%%%%%%%%%%%%%%%%%%%%%%%%%%%%%%%%%%%%%%%%%%%%%%%%%%%%%%%%%%%%%%%%%%%%%%%%%%%%%%%%%%%%%%%%%%%%%%%%%%%%%%%%%%%%%%%%%%%%
\usepackage{amsmath}

\setcounter{MaxMatrixCols}{10}
%TCIDATA{OutputFilter=LATEX.DLL}
%TCIDATA{Version=5.50.0.2960}
%TCIDATA{Codepage=1252}
%TCIDATA{<META NAME="SaveForMode" CONTENT="1">}
%TCIDATA{BibliographyScheme=Manual}
%TCIDATA{Created=Monday, September 22, 2003 09:40:11}
%TCIDATA{LastRevised=Monday, September 15, 2025 20:23:34}
%TCIDATA{<META NAME="GraphicsSave" CONTENT="32">}
%TCIDATA{<META NAME="DocumentShell" CONTENT="General\SW\Blank - Standard LaTeX Article">}
%TCIDATA{Language=American English}
%TCIDATA{CSTFile=LaTeX article (bright).cst}

\newtheorem{theorem}{Theorem}
\newtheorem{acknowledgement}[theorem]{Acknowledgement}
\newtheorem{algorithm}[theorem]{Algorithm}
\newtheorem{axiom}[theorem]{Axiom}
\newtheorem{case}[theorem]{Case}
\newtheorem{claim}[theorem]{Claim}
\newtheorem{conclusion}[theorem]{Conclusion}
\newtheorem{condition}[theorem]{Condition}
\newtheorem{conjecture}[theorem]{Conjecture}
\newtheorem{corollary}[theorem]{Corollary}
\newtheorem{criterion}[theorem]{Criterion}
\newtheorem{definition}[theorem]{Definition}
\newtheorem{example}[theorem]{Example}
\newtheorem{exercise}[theorem]{Exercise}
\newtheorem{lemma}[theorem]{Lemma}
\newtheorem{notation}[theorem]{Notation}
\newtheorem{problem}[theorem]{Problem}
\newtheorem{proposition}[theorem]{Proposition}
\newtheorem{remark}[theorem]{Remark}
\newtheorem{solution}[theorem]{Solution}
\newtheorem{summary}[theorem]{Summary}
\newenvironment{proof}[1][Proof]{\textbf{#1.} }{\ \rule{0.5em}{0.5em}}
\input{tcilatex}
\input{tcilatex}
\setlength{\textwidth}{6.5in}
\setlength{\textheight}{9in}
\setlength{\topmargin}{-.50in}
\setlength{\oddsidemargin}{.0in}
\setlength{\evensidemargin}{.0in}

\begin{document}


\begin{center}
{\LARGE Assignment 4}
\end{center}

\noindent Blankenau\hspace*{\stretch{0.7500}}

\noindent Economics 905, Fall 2025

\noindent \textit{Instructions.} You may work in groups and may compare
answers prior to submission. Joint submissions are allowed and encouraged.
All work should be legible and concise. A due date will be announced in
class.

\begin{enumerate}
\item Consider an economy where a representative, two-period-lived agent
receives $y_{1}$ units of a consumption good as an endowment in the first
period of life and none in the second. In the second period the
representative agent receives 1 unit of time which is inelastically supplied
in a competitive labor market. The agent can save some of the consumption
good while young. This becomes capital in the subsequent period which can be
rented to a firm. A representative firm can transform labor and capital into
the consumption good according to $y_{2}=zk^{\alpha }n^{1-\alpha }$ where $k$
and $n$ are capital and labor hired by the firm in period 2. \ This is a 
\textit{learning-by-doing} economy where the more that is produced economy
wide, the more productive any firm is. Since output is increasing in
capital, we can capture this by letting $z=\bar{k}^{\gamma }$ where $\bar{k}$
is the average capital stock and $\gamma >0$. Each agent and firm takes $z$
as parametric in making choices but in equilibrium $z=\bar{k}^{\gamma }.$ In
the second period, the firm chooses $k$ and $n$ to maximize profits. The
depreciation rate on capital is 1. The representative agent chooses capital
holdings $c_{1}$ and $c_{2}$ to maximize $\ln c_{1}+\ln c_{2}$ subject to
her budget constraints where $c_{1}\ $and $c_{2}$ are consumption in the
first and second periods. All markets are competitive.

\begin{enumerate}
\item Set up the social planner's problem and find the Pareto optimal
capital holdings.

\item Find the capital holdings that would arise in a competitive
equilibrium.

\item An economist has taken Econ 905 and recognizes that the competitive
equilibrium is inefficient. She proposes to subsidize the return to savings
so that the representative agent receives a return of $r\left( 1+\varsigma
\right) $ per unit of capital rented to the firm rather than $r$ where $r$
is the rental price of capital paid by the firm and $\varsigma r$ is the
subsidy paid by the government per unit of capital. This subsidy is to be
funded with a lump tax, $\tilde{\tau}$, in the second period on the
representative agent (no taxes while young). While the agent takes $\tilde{%
\tau}$ as given, this is proportional to output in equilibrium so in
equilibrium $\tilde{\tau}=\tau y.$ What is the optimal extent of this
policy? That is, how should the government set $\tau $ and $\varsigma \ $if
it wishes to maximize the utility of the representative agent?
\end{enumerate}

\item Consider a two-period economy with 2 agents. Agent 1 and agent 2 are
each endowed with one unit of the consumption good in period 1 and one unit
of the consumption good in period 2. Agents differ in preferences which are
give by%
\begin{equation*}
U_{i}\left( c_{i}\left( 1\right) ,c_{i}\left( 2\right) \right) =\phi _{i}\ln
c_{i}\left( 1\right) +\left( 1-\phi _{i}\right) \ln \left( c_{i}\left(
2\right) \right)
\end{equation*}%
where $c_{i}\left( j\right) $ is consumption in period $j$ by agent $i$ for $%
j,i\in \left\{ 1,2\right\} \ $and the preference parameter $\phi _{i}$ is
between 0 and 1. Agents can issue or purchase 1 period bonds in the first
period that yield $r$ units of the consumption good in period 2 per unit of
bond purchased in period 1. There is no default on bonds.

\begin{enumerate}
\item Find bond holdings for agent 1 as a function of $\phi _{1}\ $and $r\ $%
and find bond holdings for agent 2 as a function of $\phi _{2}\ $and $r.$
The problems are closely related so you can solve the second one by analogy.
Find the equilibrium interest rate as a function of $\phi _{1}$ and $\phi
_{2}.$

\item Now suppose there is a third agent. Agent 3 is also endowed with one
unit of the consumption good in period 1 and one unit of the consumption
good in period 2. This agent has the same sort of preferences; i.e.%
\begin{equation*}
U_{3}\left( c_{3}\left( 1\right) ,c_{3}\left( 2\right) \right) =\phi _{3}\ln
c_{i}\left( 1\right) +\left( 1-\phi _{3}\right) \ln \left( c_{i}\left(
2\right) \right) .
\end{equation*}%
Find the equilibrium interest rate.

\pagebreak
\end{enumerate}
\end{enumerate}

\begin{enumerate}
\item ...

\item \ \ \ \ \ \ 

\begin{enumerate}
\item Objective is to max%
\begin{equation*}
\phi _{i}\ln \left( 1-b_{i}\right) +\left( 1-\phi _{i}\right) \ln \left(
1+rb_{i}\right) 
\end{equation*}%
FOC is

\item 
\begin{equation*}
\frac{\phi _{i}}{1-b_{i}}=\frac{r\left( 1-\phi _{i}\right) }{1+rb_{i}}
\end{equation*}%
Use this to solve for $b_{i}$ 
\begin{eqnarray*}
\phi _{i}\left( 1+rb_{i}\right)  &=&r\left( 1-\phi _{i}\right) \left(
1-b_{i}\right)  \\
\phi _{i}+\phi _{i}rb_{i} &=&\left( 1-\phi _{i}\right) \left(
r-rb_{i}\right)  \\
\phi _{i}+\phi _{i}rb_{i} &=&r-rb_{i}-\phi _{i}r+\phi _{i}rb_{i} \\
\phi _{i} &=&r-rb_{i}-\phi _{i}r \\
rb_{i} &=&r-\phi _{i}r-\phi _{i} \\
b_{i} &=&1-\phi _{i}-\frac{\phi _{i}}{r}
\end{eqnarray*}

Bond market clearing requires 
\begin{equation*}
1-\phi _{1}-\frac{\phi _{1}}{r}+1-\phi _{2}-\frac{\phi _{2}}{r}=0
\end{equation*}%
use this to solve for $r$.%
\begin{eqnarray*}
r\left( 1-\phi _{1}\right) -\phi _{1}+\left( 1-\phi _{2}\right) r-\phi _{2}
&=&0 \\
r\left( 2-\phi _{1}-\phi _{2}\right)  &=&\phi _{1}+\phi _{2}
\end{eqnarray*}%
\begin{equation*}
r=\frac{\phi _{1}+\phi _{2}}{2-\phi _{1}-\phi _{2}}
\end{equation*}

\item We can solve this agent's problem by analogy. Then market clearing is 
\begin{equation*}
1-\phi _{1}-\frac{\phi _{1}}{r}+1-\phi _{2}-\frac{\phi _{2}}{r}+1-\phi _{3}-%
\frac{\phi _{3}}{r}=0
\end{equation*}%
or%
\begin{equation*}
3-\phi _{1}-\phi _{2}-\phi _{3}=\frac{\phi _{1}}{r}+\frac{\phi _{2}}{r}+%
\frac{\phi _{3}}{r}
\end{equation*}%
or%
\begin{equation*}
r=\frac{\phi _{1}+\phi _{2}+\phi _{3}}{3-\phi _{1}-\phi _{2}-\phi _{3}}.
\end{equation*}
\end{enumerate}
\end{enumerate}

\end{document}
